\begin{textsample}{NEG Dim 7 – human – Score -2.00 – sp/sp\_ef\_linguagens\_lp\_5\_p23.txt}  \label{ex:f7_neg_005}
Ainda decorrente da opção teórica geral do documento, que considera a língua numa perspectiva enunciativo-discursiva, cabe uma última palavra sobre a reflexão linguística-semiótica : além da continuidade do estudo da ortografia, pontuação e acentuação em suas regularidades e irregularidades, são aprofundados, progressivamente, os estudos que regem a língua dentro da norma padrão.
Mas é importante ressaltar que esses devem estar articulados aos outros \textbf{eixos} de aprendizagem.
Isso significa que o estudo da língua deve se \textbf{colocar} a favor da construção do sentido, do reconhecimento das estratégias do dizer.

% matched lemmas: colocar, eixo
\end{textsample}
